% ===========================================================================
%  Bagian 5.6 -- hb_area
% ===========================================================================
\section{Model hierarki Bayes area-level: \texttt{hb\_area}}\label{sec:hbarea}

\subsection{Setting data dan kapan dipakai}\label{sec:hbarea-setting}

Data berbentuk FH (satu baris per domain: respons $y_d$, kovariat $x_d$, dan
opsional $\psi_d$), tetapi estimasi dilakukan sepenuhnya dalam kerangka Bayes
lewat \pkg{INLA} \citep{rue2009approximate, lindgren2015bayesian}. Fungsi ini
paling luas cakupannya di paket ini: enam family likelihood, enam struktur
spasial, lima struktur temporal, dan tujuh pilihan interaksi
spatio-temporal. Cocok ketika (i) indikator non-Gaussian (rasio, proporsi,
jumlah) perlu dimodelkan langsung, atau (ii) struktur spasial/temporal
eksplisit diinginkan dengan ketidakpastian penuh posterior.

\subsection{Persamaan model}\label{sec:hbarea-model}

Kerangka umum yang dirangkai fungsi ini adalah model latent Gaussian:
\begin{equation}
  g(\mu_d)=x_d'\beta+\varphi_d,\qquad
  y_d\mid\mu_d\sim\mathcal F(\mu_d;\tau),\qquad
  \varphi\sim N\bigl(0,\;\tau^{-1}Q^{-1}\bigr),
  \label{eq:hbarea-model}
\end{equation}
dengan $g$ link family (default identitas, logit, log; binomial dan beta juga
menerima probit, cloglog, loglog lewat \code{inla\_engine(link = )}) dan $Q$ struktur
presisi yang dipilih lewat argumen \code{spatial}, \code{temporal}, dan
\code{st\_interaction}. Struktur spasial yang tersedia persis seperti yang
didefinisikan \pkg{INLA} (baris 570-649 \code{R/hb\_area.R}):
\begin{itemize}
  \item \code{"none"}: $\varphi_d$ i.i.d.\ --- hanya dibuat bila tidak ada
        interaksi lain (baris 571-574);
  \item \code{"bym2"}: BYM2 dengan presisi $Q=\tau\bigl(s\,\varphi_b Q_{\text{bym2}}+
        (1-s)\Imat\bigr)$ pada skala parsimonius (bagian~\ref{sec:bym2});
  \item \code{"bym"}: BYM klasik, dua parameter presisi
        (\code{prec.unstruct}, \code{prec.spatial});
  \item \code{"besag"}: ICAR murni $Q_{\text{ICAR}}$;
  \item \code{"generic1"} (alias \code{"leroux"}): $Q=\tau\bigl(\Imat-(\beta/\lambda_{\max})C\bigr)$
        dengan $C=\Imat-R$ dan $R$ matriks struktur ICAR (bentuk Leroux);
  \item \code{"slm"}: model lag spasial $x=(\Imat-\rho\Wmat)^{-1}
        (\Xmat\beta+\varepsilon)$ dengan $\rho$ dipetakan ke $[0,1]$ lewat
        rentang akar eigen $\Wmat$.
\end{itemize}
Struktur temporal: \code{"none"}, \code{"rw1"}, \code{"rw2"}, \code{"ar1},
\code{"iid"} pada indeks waktu \code{..time\_id..} (baris 651-728), dengan
interaksi opsional yang dipetakan ke empat tipe Knorr--Held
(Tabel~\ref{tab:kh}).

\begin{table}[htbp]
\centering
\caption{Pemetaan \texttt{st\_interaction} ke bentuk interaksi Knorr--Held.
Sumber: penyusunan string \texttt{f(...)} di \texttt{R/hb\_area.R:651-728}.}
\label{tab:kh}
\small
\begin{tabular}{@{}l l l@{}}
\toprule
Pilihan kode & Bentuk presisi & Padanan KH \\
\midrule
\code{none}          & $Q_{\text{time}}$ saja (efek utama waktu) & --- \\
\code{domain-specific} & grup per domain, model temporal & (dipakai \pkg{tipsae}) \\
\code{separable}     & $Q_{\text{space}}\otimes Q_{\text{time}}$ & --- \\
\code{type1}         & $\Imat_{DT}$ i.i.d.\ + efek waktu & KH I \\
\code{type2}         & $\Imat_D\otimes Q_{\text{time}}$ + efek utama & KH II \\
\code{type3}         & $Q_{\text{ICAR}}\otimes\Imat_T$ + efek utama & KH III \\
\code{type4}         & $Q_{\text{ICAR}}\otimes Q_{\text{time}}$ + efek utama & KH IV \\
\bottomrule
\end{tabular}
\end{table}

Interaksi KH II--IV adalah term \emph{tambahan} pada salinan indeks
(\code{..time\_id2..}, \code{..domain\_id2..}), karena \pkg{INLA} hanya
mengizinkan satu \code{f()} per kovariat. type2 memakai
\code{f(..time\_id2.., model = temporal, group = ..domain\_id..,
control.group = iid)}; type3/type4 memakai
\code{f(..domain\_id2.., model = "besag", group = ..time\_id..)} dengan
\code{control.group} \code{iid} (type3) atau \code{temporal} (type4).
type3/type4 mensyaratkan efek spasial terstruktur
(\code{bym2}/\code{bym}/\code{besag}/\code{leroux}).
\code{separable} dengan \code{leroux} memakai \code{besagproper2}, model
Leroux yang sama, karena \code{generic1} ber-group menolak diagonal
\code{Cmatrix} negatif.

\implnote{Sebelum rilis 1.0.0, type2 selalu gagal (indeks waktu ganda,
HB-1), type3/type4 dengan \code{spatial} $\ne$ \code{"none"} selalu gagal
(indeks domain ganda, HB-2), dan type3/type4 dengan
\code{spatial = "none"} diam-diam tidak membuat interaksi (HB-3). Kini
kombinasi terakhir dihentikan dengan pesan, demikian juga
\code{slm} + \code{separable} (HB-5).}

\subsection{Prediktor titik}\label{sec:hbarea-prediktor}

Prediksi diambil dari ringkasan nilai terpasang posterior (\code{inla\_utils.R:296-305}):
\begin{equation}
\begin{aligned}
  \hat\theta_d&=\mathbb E[\mu_d\mid y]\ \text{(skala respons, lewat link invers)},
  & \mathrm{sd}_d&=\sqrt{\Var(\mu_d\mid y)},\\
  \mathrm{CI}_{95\%}&=\bigl(q_{0.025},q_{0.975}\bigr)
    \ \text{(kuantil posterior)}.
\end{aligned}
  \label{eq:hbarea-pred}
\end{equation}
Untuk area-level, nilai terpasang INLA memang sudah mewakili rata-rata area,
sehingga tidak ada transformasi skala tambahan. Kolom turunan
dihitung sesuai family: \code{trials} $\to$ \code{estimated\_total}
$=\hat\theta\times\text{trials}$; \code{exposure} $\to$ \code{rate} dan
\code{estimated\_count} $=\hat\theta\times\text{exposure}$.

\implnote{Domain tanpa sampel tidak pernah membuat baris baru: baris yang
hilang dari \code{data} tidak dievaluasi INLA, sedangkan baris dengan
\code{y = NA} tetap dievaluasi karena \code{na.pass} (HB-17). Jadi
``unsampled'' di sini berarti \emph{baris ada tetapi responsnya NA}, bukan
domain yang tidak ada dalam data --- kontras dengan \fct{hb\_unit} yang
menambahkan satu baris \code{NA} per domain.}

\subsection{Estimasi komponen varians}\label{sec:hbarea-varians}

Tidak ada Fisher scoring: seluruh inferensi diserahkan ke
\code{INLA::inla()} dengan pengaturan berikut (\code{R/hb\_area.R:738-783}):
\begin{itemize}
  \item \code{control.compute} $=$ \code{dic, waic, cpo, mlik, config}
        semuanya \code{TRUE};
  \item \code{control.inla} $=$ \code{strategy} pilihan pengguna
        (\code{"laplace"} default);
  \item \code{control.predictor} $=$ \code{compute = TRUE, link = 1}, dan
        \code{control.family\$control.link\$model} $=$ link pilihan
        pengguna (divalidasi per family; sebelum 1.0.0 argumen \code{link}
        tidak pernah dipakai, HB-6). Draws posterior untuk
        \emph{self-benchmarking} dan \fct{benchmark\_sae} memakai invers-link
        yang sama;
  \item prior komponen acak spasial/temporal: \code{prior\_prec} default
        \code{pc.prec(1, 0.01)}; prior korelasi spasial \code{prior\_phi}
        default \code{pc(0.5, 0.5)}; \code{prior\_prec\_time} default juga
        \code{pc.prec(1, 0.01)} (baris 559-566);
  \item \code{control.fixed} hanya diset untuk \emph{self-benchmarking} WFQ
        (prior difus agar kovariat augmentasi tidak menyusut, baris
        765-775); selain itu berlaku default \pkg{INLA}: rata-rata tetap
        $0$, presisi tetap $0.001$ (intersep: presisi $0$) (HB-21);
  \item argumen \code{INLA::inla()} tambahan dari
        \code{inla\_engine(...)} digabung dengan \code{modifyList}, sehingga
        daftar \code{control.*} dapat ditimpa per elemen (baris 762-763;
        HB-13 tidak berlaku lagi).
\end{itemize}
Varians efek acak yang dilaporkan (\code{random\_effect\_var}) adalah rerata
posterior varians $\mathbb E[1/\tau\mid y]$ dari marginal presisinya
(\code{.post\_mean\_variance}, \code{inla\_utils.R:20-36}), bukan
$1/\mathbb E[\tau\mid y]$ yang lebih kecil karena ketaksamaan Jensen.
Untuk \code{spatial = "bym"} yang dilaporkan adalah komponen spasial
(sebelum 1.0.0: komponen iid, HB-9), dan term interaksi ruang--waktu tidak
pernah dibaca sebagai efek utama.
Likelihood Gaussian dengan \code{vardir} memakai \code{scale} $=1/\psi_d$
dan presisi tetap $1$ sehingga $\Var(y_d)=\psi_d$ persis; family lain
memakai \code{scale}/(\code{Ntrials}, \code{E}) sesuai parameterisasi
\pkg{INLA} (baris 785-841). Pengaturan engine diberikan lewat
\code{engine = inla\_engine(strategy, prior\_*, ...)} atau
\code{engine = jax\_engine(...)} (MCMC NumPyro). Argumen lama
\code{method = "laplace"} kini usang dan dipetakan ke
\code{inla\_engine(strategy = "laplace")}; jalur \code{lme4}
(HB-14, HB-15) telah dihapus.

\paragraph{Perangkat engine JAX.} \code{jax\_engine(device = ...)} memilih
CPU, GPU NVIDIA (\code{"cuda"}) atau GPU Apple (\code{"metal"}). JAX mengunci
perangkatnya pada fit pertama satu sesi R; \code{init\_jax\_environment}
(\code{inst/python/hb\_area\_numpyro.py:76}) menyetel \code{JAX\_PLATFORMS}
dan \code{jax\_default\_device}. Pada \code{"auto"}, CUDA dipakai bila ada
GPU NVIDIA, selain itu CPU (juga di macOS). Untuk model spasial, efek ICAR/BYM2
dinyatakan lewat basis eigen $D\times D$ Laplacian graf (\code{.icar\_spectral},
\code{R/hb\_area\_jax.R:27}): biaya $O(D^3)$ sekali dan $O(D^2)$ per evaluasi
gradien, sehingga waktu engine JAX tumbuh cepat dengan $D$
(bagian~\ref{sec:bench-engine}).

\implnote{jax-metal 0.1.1 tidak dapat mengompilasi NUTS NumPyro apa adanya.
Dua operasinya tidak didukung: (i) \code{popcnt}
(\code{jnp.bitwise\_count}) pada indeks checkpoint U-turn, dan (ii)
\code{lax.cond} yang semua cabangnya mengeluarkan konstanta (``Unable to
serialize MPS module''). Selain itu, loop rantai yang digabung menjadi satu
kernel tidak pernah selesai di Metal. Bila perangkatnya Metal,
\code{\_patch\_numpyro\_for\_metal} (baris 42) mengganti (i) dengan hitung bit
\emph{shift}/\emph{mask} yang identik (\code{\_popcount32}, baris 25; diuji
terhadap NumPyro untuk 0--4095) dan melewatkan keluaran setiap cabang (ii)
lewat \code{lax.optimization\_barrier}, yang tidak mengubah nilai. Rantai
dijalankan sekuensial dan per langkah dari Python (baris 452--508). Draw
Metal sama dengan draw CPU hingga galat pembulatan. Bila kompilasi tetap
gagal, fit kembali ke CPU dengan peringatan (\code{.warn\_device\_fallback},
\code{R/python\_env.R:145}).}

\subsection{Estimasi MSE}\label{sec:hbarea-mse}

MSE dihitung dari momen posterior (\code{inla\_utils.R:299-304}):
\begin{equation}
  \widehat{\mathrm{MSE}}(\hat\theta_d)=\mathrm{sd}_d^2,\qquad
  \mathrm{RSE}_d=100\,\mathrm{sd}_d/\abs{\hat\theta_d},\qquad
  \mathrm{CI}=\hat\theta_d\pm z_{0.975}\,\mathrm{sd}_d\ \text{(jalur fallback)}.
  \label{eq:hbarea-mse}
\end{equation}
Tanpa benchmarking, ketidakpastian murni dari ringkasan marginal INLA.
Dengan \code{self\_benchmark = TRUE} dan \code{benchmark\_method =
"posterior"} (default), \code{inla.posterior.sample} menarik
\code{draws} sampel prediktor linier (\code{inla\_engine(draws, seed)}),
setiap sampel di-\emph{benchmark} (rasio untuk family positif, selisih untuk
Gaussian), dan $\hat\theta$, sd, MSE, serta CI dihitung dari sampel
terbenchmark (\code{.posterior\_benchmark}, \code{inla\_utils.R:399-427});
estimasi sebelum benchmarking disimpan di \code{hb\_unbenchmarked}. Fungsi
ini membutuhkan paket \pkg{sn}; tanpa itu dipakai metode WFQ (baris
358-368).

\subsection{Mode INLA dan validasi terhadap Stan}\label{sec:hbarea-stan}
\pkg{INLA} menyediakan dua formulasi. Mode \code{"classic"} (Rue, Martino \&
Chopin 2009) memasukkan prediktor linier $\eta$ ke dalam latent field sehingga
marginalnya dihitung langsung (Laplace atau \emph{simplified Laplace}). Mode
\code{"compact"}, default \pkg{INLA} sekarang (van Niekerk, Krainski, Rustand \&
Rue 2023), mengeluarkan $\eta$ dari latent field, menghitung marginalnya sesudah
inferensi sebagai Gaussian bersyarat pada hiperparameter, dan mengoreksi rataannya
dengan koreksi Variational Bayes berperingkat rendah; mode ini dirancang untuk model
\emph{data-rich} ($n \gg$ ukuran latent field). Pada SAE level area, $n$ kira-kira
sama dengan jumlah efek area, sehingga keuntungan kecepatannya hampir tidak ada.

\code{hb\_area} kini memakai \code{inla.mode = "classic"} untuk \code{family = "gaussian"} dan mempertahankan \code{compact} untuk family lain. Dasarnya perbandingan
dengan Stan (\pkg{cmdstanr} 2.39, 4 rantai $\times$ 4000 draw, tanpa divergensi,
ESS $\ge 4700$) dengan model dan prior identik: prior PC pada presisi
($\sigma \sim \mathrm{Exp}(-\log 0{,}01)$), ICAR ter-\emph{scale} dengan kendala
jumlah nol, BYM2 parametrisasi Riebler dengan $\phi\sim U(0,1)$ (di \pkg{INLA}
\code{prior\_phi = logitbeta(1, 1)}), $\beta$ dengan prior bawaan \pkg{INLA}.
Data: grid $7\times7$ (84 sisi), 12 area tak tersampel; RW1 pada \code{mys\_panel}.
Kolom memuat selisih rata-rata maksimum dibagi SD posterior Stan dan rentang
rasio SD \pkg{INLA}/Stan atas semua area.

\begin{center}
\begin{tabular}{lcccc}
\toprule
 & \multicolumn{2}{c}{\code{compact}} & \multicolumn{2}{c}{\code{classic}} \\
Konfigurasi & $|\Delta\mu|/\mathrm{sd}$ & rasio SD & $|\Delta\mu|/\mathrm{sd}$ & rasio SD \\
\midrule
Gaussian iid (3 data set)   & $\le 0{,}24$ & $1{,}03$--$1{,}18$ & $\le 0{,}03$ & $0{,}96$--$1{,}02$ \\
Gaussian Besag (3)          & $\le 0{,}23$ & $1{,}02$--$1{,}16$ & $\le 0{,}06$ & $0{,}92$--$1{,}02$ \\
Gaussian BYM2 (3)           & $\le 0{,}25$ & $1{,}02$--$1{,}15$ & $\le 0{,}02$ & $0{,}95$--$1{,}02$ \\
Binomial iid / BYM2         & $\le 0{,}05$ & $0{,}97$--$1{,}02$ & $\le 0{,}03$ & $0{,}98$--$1{,}02$ \\
Poisson iid / BYM2          & $\le 0{,}08$ & $0{,}96$--$1{,}06$ & $\le 0{,}02$ & $0{,}96$--$1{,}02$ \\
Beta iid / BYM2             & $\le 0{,}10$ & $0{,}97$--$1{,}02$ & $\le 0{,}02$ & $0{,}97$--$1{,}02$ \\
Gamma iid / Besag           & $\le 0{,}09$ & $0{,}96$--$1{,}03$ & $\le 0{,}02$ & $0{,}97$--$1{,}02$ \\
Gaussian iid + RW1 (\code{mys\_panel}) & $0{,}04$ & $0{,}97$--$1{,}03$ & $0{,}03$ & $0{,}97$--$1{,}03$ \\
\bottomrule
\end{tabular}
\end{center}

Pada \code{mys} ($\sigma_u^2\approx1{,}9$, besar dibanding \code{vardir}) kedua mode
akurat terhadap posterior eksak (integrasi numerik 1D). Galat Monte Carlo Stan sekitar
$0{,}013$ SD untuk rata-rata dan $1\%$ untuk SD. Untuk Gaussian waktu kedua mode sama
(D sampai 2025, BYM2 $\approx 10$ detik). Untuk family lain \code{classic} jauh lebih lambat
dan tumbuh kira-kira kuadratik dalam $n$ (beta iid: $n=1000$ $9{,}9$ vs $1{,}0$ detik,
$n=2000$ $32$ vs $1{,}3$ detik; binomial $n=2000$ $12$ vs $1{,}6$ detik), padahal
\code{compact} sudah akurat (rerata $\le 0{,}16$ SD, rasio SD $0{,}95$--$1{,}08$ atas semua run);
karena itu \code{classic} hanya untuk Gaussian. Kelemahan \code{classic}: ringkasan $\sigma^2$ Gaussian
(\code{random\_effect\_var}) bias ke bawah saat $\sigma^2$ kecil (misalnya $0{,}055$ vs
Stan $0{,}118$), sedangkan \code{compact} mendekati Stan; estimasi area tetap akurat.
Belum divalidasi: Leroux/\code{generic1}, SLM, interaksi ruang--waktu tipe I--IV, dan
\code{nbinomial}.

\subsection{Pseudo-code}\label{sec:hbarea-algo}

\begin{algorithm}[htbp]
\caption{\texttt{hb\_area} --- alur sebenarnya (\texttt{R/hb\_area.R:231-469}, \texttt{.fit\_inla\_area:473-908})}
\label{alg:hbarea}
\KwIn{formula, \code{data}, \code{domain}, \code{time}, \code{family},
  \code{spatial}, \code{temporal}, \code{st\_interaction}, $W$, \code{vardir},
  \code{trials}, \code{exposure}, \code{engine}}
\KwOut{objek \code{fastsae\_hb\_area} berisi \code{df\_hb}, \code{fit}, \code{hyperpar}}
\code{match.arg} untuk argumen pilihan; \code{"leroux"} $\to$ \code{"generic1"} (255-260)\;
\code{engine} $\leftarrow$ \texttt{.resolve\_inla\_engine} (argumen lama $\to$ \code{inla\_engine}, dengan peringatan) (261-278)\;
\code{domain\_vec} $\leftarrow$ \code{.get\_variable}; \code{time} wajib bila
  \code{temporal} $\ne$ \code{"none"} (286-310)\;
\code{mf} $\leftarrow$ \texttt{model.frame(na.pass)};\quad
  $y\leftarrow$ respons (313-314)\;
\uIf{family $=$ binomial}{\code{trials} wajib;\quad
  $y\leftarrow\texttt{round}(y\cdot\text{trials})$ bila pecahan (323-339)}
\uIf{family $=$ gamma}{$y>0$ wajib (341-346)}\;
$W_{\text{obj}}\leftarrow$ \texttt{.convert\_spatial\_weights} (atau dari \code{sf}) bila
  \code{spatial} $\ne$ \code{"none"} (348-356)\;
\If{\code{self\_benchmark}}{siapkan target; WFQ: augmentasi model (358-395)}
\If{\code{engine} $=$ \code{jax\_engine()}}{\KwRet \texttt{.fit\_jax\_area} (MCMC NumPyro) (400-421)}\;
\caption*{}
\tcp*{\emph{.fit\_inla\_area} (473-908)}
indeks \code{..domain\_id..}, \code{..obs\_id..}, \code{..time\_id..} (520-538)\;
$\texttt{hyper\_time}\leftarrow\{\texttt{prec}: \texttt{pc.prec}(1,0.01)\}$; tambah
  \texttt{rho} bila \code{ar1} (559-566)\;
baris spasial $\leftarrow$ Tabel~\ref{tab:kh}/teks bagian~\ref{sec:hbarea-model} (570-649)\;
baris temporal + interaksi $\leftarrow$ Tabel~\ref{tab:kh} (651-728)\;
susun formula \code{y \textasciitilde{} ...} (\code{"slm"}: tanpa intercept) (730-736)\;
\code{control.compute} (dic, waic, cpo, mlik, config), \code{control.predictor}
  (\code{link = 1}), \code{control.inla}; gabung argumen \code{inla\_engine} (738-783)\;
\For{family}{sesuaikan \code{scale}, \code{Ntrials}, \code{E},
  \code{control.family} --- \eqref{eq:hbarea-model} (785-841)}\;
\code{fit} $\leftarrow$ \texttt{tryCatch(do.call(INLA::inla, ...))} (843-852)\;
$\hat\theta,\mathrm{sd},\mathrm{CI},\texttt{mse},\texttt{rse}\leftarrow$
  \eqref{eq:hbarea-pred}--\eqref{eq:hbarea-mse} (\code{.extract\_inla\_results}, 854-875)\;
\If{\code{self\_benchmark} posterior}{\texttt{.posterior\_benchmark} (877-899)}
perakitan \code{df\_hb}, \code{hyperpar}, \code{goodness}
  (DIC, WAIC, marginal likelihood), kelas \code{c("fastsae\_hb\_area","fastsae\_hb","fastsae")}\;
\end{algorithm}

\subsection{Signature, argumen, objek keluaran, dan contoh}\label{sec:hbarea-api}

\begin{lstlisting}[language=R,caption={Signature \texttt{hb\_area} (\texttt{R/hb\_area.R:231-253}) dan \texttt{inla\_engine} (\texttt{R/engine.R:83}).}]
hb_area(formula, data, domain = NULL, time = NULL,
        family = c("gaussian","binomial","poisson","nbinomial","beta","gamma"),
        spatial  = c("none","bym2","bym","besag","generic1","leroux","slm"),
        temporal = c("none","rw1","rw2","ar1","iid"),
        st_interaction = c("none","domain-specific","separable",
                           "type1","type2","type3","type4"),
        W = NULL, vardir = NULL, trials = NULL, exposure = NULL,
        engine = "inla",   # atau inla_engine(...) / jax_engine(...)
        self_benchmark = FALSE, benchmark_weight = NULL,
        benchmark_target = NULL, benchmark_group = NULL,
        benchmark_method = c("posterior","wfq"),
        print_result = TRUE, method = NULL, ...)

inla_engine(strategy = c("laplace","simplified.laplace","gaussian"),
            link = NULL, scale_model = TRUE,
            prior_prec = NULL, prior_phi = NULL, prior_rho = NULL,
            prior_prec_time = NULL, prior_rho_time = NULL,
            prior_prec_subarea = NULL, draws = NULL, seed = 0L, ...)
# NULL = default fungsi: hb_area memakai pc.prec(1, 0.01),
# pc(0.5, 0.5) untuk phi, dan draws = 1000

jax_engine(warmup = 500, samples = 1000, chains = 2,
           device = c("auto","cpu","cuda","metal"), seed = 42,
           prior_beta = c("normal","horseshoe"), robust = FALSE,
           keep_draws = TRUE, progress = interactive())
\end{lstlisting}

\begin{itemize}
  \item \code{family} menentukan kolom pendukung yang wajib: Gaussian
        $\to$ \code{vardir} (opsional; bila ada, dipakai sebagai
        \code{scale}); Binomial $\to$ \code{trials}; Poisson/Negatif
        Binomial $\to$ \code{exposure}; Beta $\to$ \code{vardir} atau
        \code{trials} untuk presisi.
  \item \code{W} --- matriks ketetanggaan (dapat berupa \code{listw},
        \code{nb}, \code{matrix}, atau \code{inla.graph}); bobot dibakar ke
        graf biner simetris (bagian~\ref{sec:bayes-w}).
  \item \code{engine} --- \code{"inla"}/\code{inla\_engine(...)} (default)
        atau \code{jax\_engine(...)}; argumen INLA yang dulu diberikan
        langsung ke \fct{hb\_area} (\code{method}, \code{strategy},
        \code{prior\_*}, \code{...}) masih diterima dengan peringatan dan
        diteruskan ke \fct{inla\_engine}.
  \item \code{W} kosong dengan \code{data} berkelas \code{sf}: matriks
        ketetanggaan queen biner dibangun otomatis lewat
        \fct{create\_weights}.
\end{itemize}

Objek keluaran: \code{c("fastsae\_hb\_area", "fastsae\_hb", "fastsae")} berisi
\code{df\_hb} (kolom \code{domain, y, hb, sd, mse, rse, ci\_lower,
ci\_upper, linear\_pred} + kolom family), \code{fit} (objek \code{inla}),
\code{hyperpar}, \code{goodness}, \code{random\_effect\_var},
\code{engine}, \code{engine\_spec}, \code{W}, \code{n\_iter},
\code{convergence}, \code{formula}, \code{method}, \code{call},
\code{data}.

\begin{lstlisting}[language=R,caption={Contoh yang dapat dijalankan (dari \texttt{man/hb\_area.Rd}).}]
library(fastsae)
if (requireNamespace("INLA", quietly = TRUE)) {
  # Gaussian Fay-Herriot tanpa spasial
  m_norm <- hb_area(y ~ x1 + x2 + x3, data = mys, vardir = "vardir",
                    family = "gaussian")

  # Gaussian dengan efek spasial BYM2
  m_bym2 <- hb_area(y ~ x1 + x2 + x3, data = mys, vardir = "vardir",
                    family = "gaussian", spatial = "bym2",
                    W = mys_proxmat)
  head(m_norm$df_hb)
}
\end{lstlisting}

\subsection{Referensi kunci}\label{sec:hbarea-ref}

\citet{rue2009approximate} untuk metode INLA, \citet{lindgren2015bayesian}
untuk implementasi \pkg{R-INLA}, \citet{riebler2016intuitive} untuk BYM2
(blok roxygen fungsi ini memuat lima rujukan: \citet{rao2015small},
\citet{riebler2016intuitive}, \citet{marhuenda2013small},
\citet{denicolo2024tipsae}, \citet{rue2009approximate}; lihat blok
\code{@references} di \code{R/hb\_area.R}).
